% Documento principal de la entrega. Genera SemArq2_TeamBales.pdf:
%   cd latex && latexmk -pdf -jobname=SemArq2_TeamBales main.tex

\documentclass[11pt,a4paper]{article}
% Preámbulo común de SemArq2_TeamBales.pdf. Lo carga main.tex; los fragmentos de
% cada inciso no deben tener preámbulo propio ni cargar paquetes.

\usepackage[T1]{fontenc}
\usepackage[utf8]{inputenc}
\usepackage[spanish,es-tabla,es-nodecimaldot]{babel}
\usepackage{lmodern}
\usepackage{microtype}
\usepackage[a4paper,margin=2.5cm]{geometry}

\usepackage{graphicx}
% Las rutas de las imágenes son relativas a la raíz del repo, por ejemplo:
% \includegraphics{preguntas/p2-elasticidad-consistencia/diagramas/mi-diagrama.png}
\graphicspath{{../}}

\usepackage{booktabs}
\usepackage{tabularx}
\usepackage{float}
\usepackage{caption}
\usepackage{enumitem}
\usepackage{amsmath}
\usepackage[table]{xcolor}
\usepackage{tikz}

% ---------- Colores de la portada ----------
\definecolor{portadaoscuro}{HTML}{1E3A6E}
\definecolor{portadaacento}{HTML}{3D5FB0}
\definecolor{portadaclaro}{HTML}{C9D8F5}
\definecolor{portadagris}{HTML}{6B7280}
\definecolor{portadafila}{HTML}{EEF3FC}
\usepackage{listings}
\usepackage[hidelinks]{hyperref}

\definecolor{codigofondo}{HTML}{F6F8FA}
\definecolor{codigoclave}{HTML}{0033B3}
\definecolor{codigocomentario}{HTML}{6A737D}
\definecolor{codigocadena}{HTML}{067D17}

\lstdefinestyle{codigogeneral}{
  basicstyle=\ttfamily\small,
  keywordstyle=\color{codigoclave}\bfseries,
  commentstyle=\color{codigocomentario}\itshape,
  stringstyle=\color{codigocadena},
  backgroundcolor=\color{codigofondo},
  frame=single, rulecolor=\color{black!15},
  numbers=left, numberstyle=\tiny\color{black!50},
  breaklines=true, showstringspaces=false, tabsize=4,
  literate={á}{{\'a}}1 {é}{{\'e}}1 {í}{{\'i}}1 {ó}{{\'o}}1 {ú}{{\'u}}1
           {Á}{{\'A}}1 {É}{{\'E}}1 {Í}{{\'I}}1 {Ó}{{\'O}}1 {Ú}{{\'U}}1
           {ñ}{{\~n}}1 {Ñ}{{\~N}}1 {ü}{{\"u}}1 {¿}{{?`}}1 {¡}{{!`}}1
           {→}{{$\rightarrow$}}2 {×}{{$\times$}}1
}
\lstset{style=codigogeneral}

% ---------- Incisos ----------
% \inciso{título}{ruta del fragmento desde la raíz del repo}
% main.tex pone el título de la subsección; el fragmento solo trae el contenido.
% Si el fragmento todavía no existe, se muestra "Pendiente" y el PDF sigue compilando.
\newcommand{\inciso}[2]{%
  \subsection{#1}%
  \InputIfFileExists{../#2}{}{\textit{Pendiente.}}%
}


\begin{document}

% ---------- Portada ----------
% La tabla de integrantes debe coincidir con la tabla de miembros del README principal.
\begin{titlepage}
  \begin{tikzpicture}[remember picture, overlay]
    \fill[portadaoscuro] (current page.north west) rectangle ([yshift=-8cm]current page.north east);
    \fill[portadaacento] ([yshift=-8cm]current page.north west) rectangle ([yshift=-8.25cm]current page.north east);
    \node[anchor=north west, text=portadaclaro, font=\large\scshape]
      at ([xshift=2.5cm, yshift=-2.2cm]current page.north west) {Ingeniería de Software};
    \node[anchor=north west, text=white, font=\fontsize{26}{30}\selectfont\bfseries]
      at ([xshift=2.5cm, yshift=-3.1cm]current page.north west) {Seminario sobre Arquitectura 2};
    \node[anchor=north west, text=portadaclaro, font=\Large]
      at ([xshift=2.5cm, yshift=-5cm]current page.north west) {Arquitecturas de Software Distribuidas};
    \node[anchor=north west, text=white, font=\large]
      at ([xshift=2.5cm, yshift=-6.2cm]current page.north west) {Service-Based, Event-Driven, Space-Based, SOA y Microservicios};
    \fill[portadaoscuro] (current page.south west) rectangle ([yshift=1cm]current page.south east);
  \end{tikzpicture}

  \vspace*{6.6cm}

  \begin{center}
    {\color{portadagris}\large\scshape Equipo 7\par}
    \vspace{0.1cm}
    {\color{portadaoscuro}\fontsize{50}{56}\selectfont\bfseries Bales\par}
    \vspace{0.35cm}
    {\color{portadaacento}\rule{3.5cm}{2.5pt}\par}
  \end{center}

  \vspace{1.4cm}

  \noindent{\color{portadagris}\scshape Integrantes\par}
  \vspace{0.3cm}
  {\renewcommand{\arraystretch}{1.7}%
  \rowcolors{2}{portadafila}{white}\noindent%
  \begin{tabularx}{\textwidth}{>{\raggedright\arraybackslash}X >{\ttfamily\small}l >{\raggedright\arraybackslash}X}
    \rowcolor{portadaoscuro}
    \textcolor{white}{\textbf{Integrante}} & \textcolor{white}{\rmfamily\normalsize\textbf{GitHub}} & \textcolor{white}{\textbf{Pregunta(s)}} \\
    Heily Rodríguez Rodríguez & @heilyrodriguez225 & P1, P4 \\
     & @Mackandal04 & P5 \\
     & @QuiquiMatCom2004 & P3, P6 \\
     & @solerch5510 & P2 \\
  \end{tabularx}\par}

  \vfill

  \begin{center}
    {\color{portadagris}\small\scshape Repositorio\par}
    \vspace{0.15cm}
    {\color{portadaacento}\url{https://github.com/Team-Bales/seminario-arq-distribuidas}\par}
  \end{center}
  \vspace{0.6cm}
\end{titlepage}

\tableofcontents
\newpage

\section{P1: Análisis Comparativo de Comunicación y Acoplamiento}
\inciso{a) Acoplamiento dinámico}{preguntas/p1-comunicacion-acoplamiento/latex/inciso-a.tex}
\inciso{b) Tolerancia a fallos y consistencia}{preguntas/p1-comunicacion-acoplamiento/latex/inciso-b.tex}

\newpage
\section{P2: Diseño para Elasticidad Extrema y Consistencia de Datos}
\inciso{a) Space-Based vs. Microservices}{preguntas/p2-elasticidad-consistencia/latex/inciso-a.tex}
\inciso{b) Colisiones de datos en cachés replicadas}{preguntas/p2-elasticidad-consistencia/latex/inciso-b.tex}

\newpage
\section{P3: Análisis de Fallos y Resiliencia en Arquitecturas Dirigidas por Eventos}
\inciso{a) Impacto del fallo}{preguntas/p3-fallos-resiliencia-eda/latex/inciso-a.tex}
\inciso{b) Workflow Event Pattern}{preguntas/p3-fallos-resiliencia-eda/latex/inciso-b.tex}

\newpage
\section{P4: El Dilema de la Granularidad y el Acoplamiento en Microservicios}
\inciso{a) Error fundamental de diseño}{preguntas/p4-granularidad-microservicios/latex/inciso-a.tex}
\inciso{b) Nueva estrategia de descomposición}{preguntas/p4-granularidad-microservicios/latex/inciso-b.tex}

\newpage
\section{P5: Comparación de Atributos de Calidad y Trade-offs}
\inciso{a) Consistencia de datos}{preguntas/p5-atributos-calidad-tradeoffs/latex/inciso-a.tex}
\inciso{b) Despliegue independiente y trade-offs}{preguntas/p5-atributos-calidad-tradeoffs/latex/inciso-b.tex}

\newpage
\section{P6: El Legado de SOA y su Influencia en Arquitecturas Modernas}
\inciso{a) Reutilización en Microservices}{preguntas/p6-legado-soa/latex/inciso-a.tex}
\inciso{b) Cuándo todavía conviene orquestar}{preguntas/p6-legado-soa/latex/inciso-b.tex}

\end{document}
